\documentclass[11pt]{article}
\usepackage[margin=1in]{geometry}
\usepackage{amsmath,amssymb,amsthm,mathtools,hyperref,enumitem}
\hypersetup{hidelinks}
\newtheorem{theorem}{Theorem}
\newtheorem{lemma}[theorem]{Lemma}
\newtheorem{proposition}[theorem]{Proposition}
\newtheorem{corollary}[theorem]{Corollary}
\theoremstyle{remark}\newtheorem{remark}[theorem]{Remark}
\newcommand{\Cp}{\mathbb C_p}\newcommand{\OO}{\mathcal O_{\mathbb C_p}}
\newcommand{\pos}[1]{\left(#1\right)_+}
\title{A polynomial growth exponent for finite families\protect\\in torsion specialization}
\author{Research note}
\date{27 November 2026}
\begin{document}\maketitle
\begin{abstract}
Fix the number $n$ of integral $p$-adic polynomials and a positive tolerance $\delta$.
We prove that the worst necessary torsion-order exponent, at total valuation budget $B$,
is $B^{n+o(1)}$. More precisely, it lies between a positive constant times $B^n$ and a
constant times $B^n(\log B)^2\log\log B$. The result determines the growth exponent
asked for in Question \texttt{cost:q:cardinality(b)} of the pinned ProveIt article~\cite{proveit-cost}.
The proof combines root accounting with a shifted product potential and an explicit
construction on consecutive blocks of torsion shells. It is unconditional and does
not use RH or GRH. It does not determine a leading asymptotic constant.
\end{abstract}

\section{Precise setting and result}
Fix a prime $p$, an integer $h_0\ge1$, and $\rho\ge r_{h_0}$, where
\[
 e_h=(p-1)p^{h-1},\qquad r_h=e_h^{-1},\qquad
 Z_h=\{\zeta-1:\operatorname{ord}(\zeta)=p^h\}.
\]
Write $v(p)=1$ and $v(0)=+\infty$. The shell $Z_h$ contains $e_h$ points, all of
valuation $r_h$. For a nonzero polynomial $g=\sum a_jT^j\in\OO[T]$ put
\[
 G_u(g)=\min_j\{v(a_j)+ju\}.
\]
A packet of at most $n$ members consists of such polynomials $g_i$ and finite thresholds
$t_i\ge G_\rho(g_i)$, with cost $\sum_i t_i$. Its loss at $z$ is
\[
 \mathcal L(z)=\sum_i\pos{v(g_i(z))-t_i}.
\]
Here $x_+=\max(0,x)$, including $+\infty$. A packet blocks the shells through $H$ at
tolerance $\delta$ if $\mathcal L(z)>\delta$ at every point of
$Z_{h_0}\cup\cdots\cup Z_H$.
Let $H_n(B,\delta)$ be the least $H\ge h_0$ for which every packet of at most $n$ members
and cost at most $B$ has a point of loss at most $\delta$ among those shells.

\begin{theorem}[Positive-tolerance growth exponent]\label{thm:main}
Fix $p,h_0,n$ and $\delta>0$, with $n\ge1$. There are positive constants $c,C,B_0$,
depending only on these parameters, such that for $B\ge B_0$,
\[
 cB^n\le H_n(B,\delta)
 \le C B^n(\log B)^2\log\log B.
\]
Consequently
\[
 \lim_{B\to\infty}\frac{\log H_n(B,\delta)}{\log B}=n.
\]
Both bounds are uniform in $\rho\ge r_{h_0}$, the degrees of the polynomials, and
the ramification of their coefficients. The lower examples have coefficients in
finite cyclotomic extensions and thresholds $t_i=G_\rho(g_i)=v(g_i(0))$.
For $n=1$ the sharper order estimate $H_1(B,\delta)=\Theta(B)$ holds.
\end{theorem}

This is a statement about the exponent $h$ in the torsion order $p^h$.
It gives $H_n(B,\delta)=B^{n+o(1)}$, not a polynomial bound on the torsion order itself.
It also keeps $\delta>0$ fixed. The constants are not uniform as $\delta\downarrow0$.

\section{Reduction to roots on the tested shells}
For completeness we justify the root model used in the proof. If
$g=b\prod_\alpha(T-\alpha)$, with multiplicities, then
\begin{equation}\label{eq:gauss}
 G_u(g)=v(b)+\sum_\alpha\min\{u,v(\alpha)\}.
\end{equation}
This follows from multiplicativity of the Gauss valuation. In particular
$G_\rho(g)-G_0(g)$ is the sum of the nonnegative quantities
$\min(\rho,v(\alpha))-\min(0,v(\alpha))$.

For a finite set of tested shells, one may replace each root with valuation $r_h$
by a nearest point $\beta\in Z_h$. The ultrametric inequality gives
$v(z-\alpha)\le v(z-\beta)$ for every $z\in Z_h$; on other shells the contribution
depends only on the valuation of the root. Remove all roots whose valuations do not
belong to the tested shells, take leading coefficient one, and reset the threshold
to the new Gauss valuation. A removed root had a nonpositive contribution to
$v(g(z))-G_\rho(g)$ at every test point, so no positive loss decreases.
The new cost is at most the old cost by \eqref{eq:gauss} and $G_0(g)\ge0$.
If a member has more than $e_h$ roots on shell $h$, replace them by every point of
$Z_h$, once each. Its loss on this shell becomes infinite; on other shells, removing
roots with this valuation cannot decrease its excess. Thus the packet still blocks.

It therefore suffices to study packets with monic polynomials whose roots are on the
tested shells and with at most $e_h$ roots of each member on each shell. Pad with
constant polynomials to obtain exactly $n$ members. Write
\begin{align*}
 m_{i,h}&=\text{number of roots of member $i$ on shell $h$},&
 w_{i,h}&=m_{i,h}r_h,\\
 W_i(h)&=\sum_{k=h_0}^{h}w_{i,k},&
 S_i(h)&=\sum_{h_0\le k<h}w_{i,k}(1-p^{k-h}),\\
 E_{i,h}(z)&=\sum_{\alpha\text{ on shell }h}
                 \bigl(v(z-\alpha)-r_h\bigr).&&
\end{align*}
Then $0\le w_{i,h}\le1$, $\sum_i W_i(H)\le B$, and the excess of member $i$ at
$z\in Z_h$ is exactly
\begin{equation}\label{eq:excess}
 v(g_i(z))-G_\rho(g_i)=E_{i,h}(z)-S_i(h).
\end{equation}
Set $W_i(h_0-1)=0$. These reductions also show that for a finite $H$ the minimum
blocking cost is attained: there are only finitely many reduced root multisets.

\section{A capped average with positive tolerance}
\begin{lemma}[One-root average]\label{lem:average}
For $\alpha\in Z_h$ and $c>0$,
\[
 \frac1{e_h}\sum_{z\in Z_h}\min\{v(z-\alpha)-r_h,c\}
 \le (h+c)r_h.
\]
\end{lemma}
\begin{proof}
Let $P_h(T)=\Phi_{p^h}(1+T)$. Differentiating the cyclotomic quotient at $\alpha$ gives
$v(P_h'(\alpha))=h-1/(p-1)$. Thus the sum over $z\ne\alpha$ of
$v(z-\alpha)-r_h$ equals
$h-1/(p-1)-(e_h-1)r_h\le h$. Capping can only decrease these terms; the term at
$z=\alpha$ contributes $c$. Divide by $e_h$.
\end{proof}

\begin{lemma}[Forced shifted growth]\label{lem:forced}
On every blocked shell $h$, with $w_h=\sum_i w_{i,h}$, one has
\begin{equation}\label{eq:forced}
 1\le h\sum_{i=1}^n\frac{w_{i,h}}{S_i(h)+\delta}+w_h.
\end{equation}
\end{lemma}
\begin{proof}
For each member cap each root's excess at $c_i=S_i(h)+\delta$ and denote the sum
of capped root excesses by $\widetilde E_i(z)$. We claim that
$\sum_i\widetilde E_i(z)/c_i\ge1$ at each $z\in Z_h$.
If any root is capped, its capped value already equals $c_i$, proving the claim.
Otherwise $\widetilde E_i=E_{i,h}$. If the sum of their ratios were less than one,
each $E_{i,h}<S_i+\delta$, and
\[
 \pos{E_{i,h}-S_i}\le
   \frac{\delta E_{i,h}}{S_i+\delta}.
\]
The sum of losses would be less than $\delta$, contradicting blocking and
\eqref{eq:excess}. Average the claimed ratio inequality over $Z_h$ and apply
Lemma~\ref{lem:average} to all roots. The resulting right side is precisely
$h\sum_i w_{i,h}/c_i+\sum_i w_{i,h}$.
\end{proof}

\section{A quantitative bound from a shifted product}
\begin{proposition}[Explicit bound for every blocker]\label{prop:explicit}
Choose an integer $d\ge1$ and a real number $K>1$, and put
\[
 \lambda=\frac{1-p^{-d}}K,\qquad
 L=n\log\left(1+\frac{B}{n\delta}\right),\qquad
 D=d+\left\lceil\frac{dL}{\log K}\right\rceil.
\]
If a packet of at most $n$ members and cost at most $B$ blocks every shell from
$h_0$ through $H$ at tolerance $\delta>0$, then
\begin{equation}\label{eq:explicit}
 \left(\frac{H+1}{h_0}\right)^\lambda
 \le\left(1+\frac{B}{n\delta}\right)^n
 \left[1+\mathrm e D+
       \mathrm e(\delta+n)\log\left(1+\frac B\delta\right)\right].
\end{equation}
Here $\mathrm e=\exp(1)$.
\end{proposition}
\begin{proof}
Apply the root reduction and set
\[
 Q_i(h)=\delta+W_i(h),\qquad P_h=\prod_{i=1}^nQ_i(h),\qquad
 P_{h_0-1}=\delta^n.
\]
All factors are positive and nondecreasing. By the arithmetic--geometric mean inequality,
\begin{equation}\label{eq:Pbound}
 P_H\le\left(\delta+\frac Bn\right)^n,
 \qquad \log\frac{P_H}{P_{h_0-1}}\le L.
\end{equation}
Declare the first $d$ shells exceptional. A later shell $h$ is exceptional if
$P_{h-1}/P_{h-d}>K$. Each logarithmic increment of $P$ appears in at most $d$
such quotients, so their number is at most $dL/\log K$. Thus at most $D$ shells
are exceptional. On every other shell, all the individual ratios
$Q_i(h-1)/Q_i(h-d)$ are at most $K$, since their product is at most $K$ and
each is at least one. Keeping in $S_i$ just the roots from shells at most $h-d$ gives
\begin{equation}\label{eq:regular}
 S_i(h)+\delta
 \ge (1-p^{-d})\bigl(W_i(h-d)+\delta\bigr)
 \ge\lambda Q_i(h-1).
\end{equation}
Let $R_h=P_h/P_{h-1}$. Expanding the product of nonnegative increments,
then using \eqref{eq:regular} and Lemma~\ref{lem:forced}, gives on a nonexceptional shell
\begin{equation}\label{eq:R}
 R_h=\prod_i\left(1+\frac{w_{i,h}}{Q_i(h-1)}\right)
 \ge1+\frac{\lambda(1-w_h)_+}{h}.
\end{equation}

Introduce the comparison sequence
\[
 A_{h_0-1}=1,\qquad
 A_h=\prod_{j=h_0}^{h}\left(1+\frac\lambda j\right),\qquad F_h=\frac{A_h}{P_h}.
\]
On a nonexceptional shell, \eqref{eq:R} implies
\begin{align*}
 F_h-F_{h-1}
 &\le F_{h-1}\left(
 \frac{1+\lambda/h}{1+\lambda(1-w_h)_+/h}-1\right)\\
 &\le F_{h-1}\frac{\lambda w_h}{h}.
\end{align*}
On an exceptional shell use $R_h\ge1$ to obtain
$F_h-F_{h-1}\le F_{h-1}\lambda/h$.
Since $0<\lambda<1$, the harmonic sum bound gives
$A_{h-1}\le\mathrm e h^\lambda$, hence
$\lambda A_{h-1}/h\le\mathrm e$ for $h\ge1$.
Consequently, after summing the increments,
\begin{equation}\label{eq:Ftelescope}
 F_H\le\delta^{-n}+\mathrm e D\delta^{-n}
       +\mathrm e\sum_{h=h_0}^H\frac{w_h}{P_{h-1}}.
\end{equation}
Write $W(h)=\sum_iW_i(h)$. Expansion of the product gives
$P_{h-1}\ge\delta^{n-1}(\delta+W(h-1))$.
Also $0\le w_h\le n$. For
$x_h=w_h/(\delta+W(h-1))\le n/\delta$, the inequality
$\log(1+x)\ge x/(1+x)$ yields
\begin{align*}
 \sum_h\frac{w_h}{P_{h-1}}
 &\le\delta^{1-n}\left(1+\frac n\delta\right)
        \sum_h\log\frac{\delta+W(h)}{\delta+W(h-1)}\\
 &\le\delta^{-n}(\delta+n)\log\left(1+\frac B\delta\right).
\end{align*}
Finally, concavity gives
$\log(1+\lambda/j)\ge\lambda\log(1+1/j)$, so
\[
 A_H\ge\left(\frac{H+1}{h_0}\right)^\lambda.
\]
Combine this with $A_H=F_HP_H$, \eqref{eq:Pbound}, and \eqref{eq:Ftelescope}.
\end{proof}

\begin{proof}[Upper bound in Theorem~\ref{thm:main}]
For large $B$ put
$q_B=\log(2+B/(n\delta))$, choose
$d=\lceil\log_p q_B\rceil$ and $K=1+q_B^{-1}$.
Then $p^{-d}\le q_B^{-1}$ and
\[
 \lambda^{-1}=1+O(q_B^{-1}),\qquad
 D=O_{p,n}(q_B^2\log q_B).
\]
Taking the $\lambda$th root in \eqref{eq:explicit}, the main factor satisfies
$(1+B/(n\delta))^{n/\lambda}=O_{n,\delta}(B^n)$.
The bracket to the power $1/\lambda$ is
$O_{p,n,\delta}(q_B^2\log q_B)$: the extra power
$O(1/q_B)$ changes this by a bounded factor.
Every blocking endpoint $H$ therefore satisfies the asserted upper bound, and
increasing this bound by one bounds $H_n(B,\delta)$.
\end{proof}

\section{Explicit blockers and the matching exponent}
For $0\le m<e_h$ define
\[
 U_h(m)=\frac1{e_h}\sum_{j=1}^{h-1}e_j\left\lfloor\frac m{e_j}\right\rfloor,
 \qquad U_h(e_h)=+\infty.
\]
If $\xi_h$ is primitive of order $p^h$, take as roots $\xi_h^b-1$ for the first $m$
positive integers $b$ prime to $p$. At every unmarked leaf, the first $m$ unit residues
have at least $\lfloor m/e_j\rfloor$ representatives in its class modulo $p^j$.
The cyclotomic distance formula
\[
 v(\xi_h^a-\xi_h^b)-r_h
 =\frac{p^{v_p(a-b)}-1}{e_h}
 =\frac1{e_h}\sum_{j=1}^{v_p(a-b)}e_j
\]
therefore gives an on-shell excess of at least $U_h(m)$ at every leaf; marked leaves
have infinite excess.

\begin{lemma}[An inexpensive blocker for a late block]\label{lem:block}
Let $a\ge2$, $b\ge a$, and $\delta>0$, and suppose
$c=(\delta+3)/(a-1)<1$. There is a monic polynomial, with all roots on the shells
$a,\ldots,b$, that has excess greater than $\delta$ at every point of those shells
and cost at most
\[
 c(b-a+1)+1.
\]
The threshold can be its Gauss valuation at any radius at least $r_a$.
\end{lemma}
\begin{proof}
On each shell $h$ use $m_h=\lceil c e_h\rceil\le e_h$ of the greedy roots above.
If $m_h=e_h$, the excess is infinite at all points of $Z_h$.
Otherwise $m_h/e_h\ge c$ and
\[
 U_h(m_h)\ge (h-1)c-\frac1{p-1}\ge(h-1)c-1.
\]
Indeed, replace every floor by its argument minus one and use
$\sum_{j=1}^{h-1}e_j/e_h<1/(p-1)$.
The earlier-shell charge is at most
\[
 S(h)\le\sum_{k=a}^{h-1}\frac{m_k}{e_k}
 \le c(h-a)+\sum_{k=a}^{\infty}r_k
 \le c(h-a)+1,
\]
because $a\ge2$ and
$\sum_{k=a}^{\infty}r_k=p^{2-a}/(p-1)^2\le1$.
The excess is consequently at least
$(a-1)c-2=\delta+1>\delta$.
The same geometric sum bounds the total cost by $c(b-a+1)+1$.
\end{proof}

\begin{proof}[Lower bound in Theorem~\ref{thm:main}]
Put $Q=\delta+3$ and assume $B\ge8n(\delta+4)$.
Start at $a_1=h_0$ and let $\ell_1=\lfloor B/n\rfloor$ and $a_2=a_1+\ell_1$.
The first member is
$\prod_{h=a_1}^{a_2-1}\Phi_{p^h}(1+T)$, with threshold its Gauss valuation.
It costs exactly $\ell_1\le B/n$ and vanishes at every tested point in its block.
For $2\le i\le n$ define
\[
 \ell_i=\left\lfloor\frac{(a_i-1)B}{2nQ}\right\rfloor,
 \qquad a_{i+1}=a_i+\ell_i,
\]
and use Lemma~\ref{lem:block} on the $i$th block. Its cost is at most
$B/(2n)+1\le B/n$. Our lower bound on $B$ ensures $a_i\ge2$, $Q/(a_i-1)<1$,
and $(a_i-1)B/(2nQ)\ge2$ at every such step. Thus all members together cost at most $B$
and block every shell from $h_0$ through $a_{n+1}-1$.
Moreover,
\[
 a_2-1\ge\frac{B}{2n},\qquad
 a_{i+1}-1\ge\frac{B}{4nQ}(a_i-1)\quad(2\le i\le n).
\]
It follows that
\begin{equation}\label{eq:lower}
 H_n(B,\delta)\ge a_{n+1}
 \ge\frac{B^n}{2n\,[4n(\delta+3)]^{n-1}}.
\end{equation}
All roots used are torsion points; finitely many suffice for each $B$, so the coefficients
belong to a finite cyclotomic extension. Their nonzero roots have valuations at most $\rho$,
which gives $G_\rho(g_i)=v(g_i(0))$.

For $n=1$, the same full-shell construction gives $H_1(B,\delta)\ge h_0+\lfloor B\rfloor$.
Monotonicity in the tolerance and the repository's exact zero-tolerance formula~\cite[Theorem \texttt{cost:thm:cardinality}]{proveit-cost}
$H_1(B,0)=h_0+\lfloor Bp^{h_0-1}\rfloor$ give $H_1(B,\delta)=\Theta(B)$.
\end{proof}

\begin{corollary}[Minimum blocking cost]
Let $C_n(H,\delta)$ be the minimum cost of a packet of at most $n$ members blocking
the shells through $H$. For fixed $p,h_0,n,\delta>0$,
\[
 C_n(H,\delta)=H^{1/n+o(1)}\qquad(H\to\infty).
\]
\end{corollary}
\begin{proof}
Finite attainment from the root reduction gives
$C_n(H,\delta)\le B$ if and only if $H<H_n(B,\delta)$.
Invert Theorem~\ref{thm:main}, or bound its logarithmic factors by $B^\varepsilon$.
\end{proof}

\section{Relation to the source and remaining questions}
The immediate source is the ProveIt article~\cite{proveit-cost}, pinned at commit
\begin{quote}\small\texttt{b82291c1f75389938c7464ef161907d7d5e75b2a}.\end{quote}
We use \texttt{Article/sections/cost.tex}, particularly its root model, fixed-cardinality
tower theorem, and Question \texttt{cost:q:cardinality(b)}.
The original transferable ingredient is the simultaneous roots-of-unity selector in
\begin{quote}\small\raggedright
\path{vendor/grh-lean/grh/src/New/Congruences/Selmer/RootsOfUnity.lean}~\cite{grh-roots}.
\end{quote}
The new step here is the shift by $\delta$ in the product potential and the normalized
increment argument. The shift prevents the activation denominators from becoming
arbitrarily small. This accounts for the different growth regime.

The article's zero-tolerance theorem gives, for fixed $n\ge2$,
\[
 \log_p^{\circ(n-1)}H_n(B,0)
  =Bp^{h_0-1}+O(\log B).
\]
With unrestricted cardinality and fixed positive $\delta$, the article gives
$\log H_*(B,\delta)=B/\delta+O(1)$.
Theorem~\ref{thm:main} supplies the fixed-cardinality, positive-tolerance growth exponent.
It leaves open whether $H_n(B,\delta)=\Theta(B^n)$ for $n\ge2$, whether a leading constant
exists, and the joint transition in which $\delta$ tends to zero with $B$.

This result resolves a research question posed in this repository at the level of its
growth exponent. It is not a resolution of a prominent named conjecture. Searches of the
primary literature on torsion specialization and horizontal $p$-adic $L$-functions found
relevant nonvanishing and specialization machinery but no matching fixed-packet growth law;
that search does not establish priority.


\section{Reproducibility and verification scope}
The accompanying standard-library Python
\href{../code/verify_torsion_phase_transition.py}{verification program}
checks 267 greedy shell-capacity instances, 52 capped-root averages, 176 explicit
late-block constructions, and all 16,384 two-member subset-root packets on the first
three shells for $p=2$ at tolerance $1/4$. All arithmetic is exact. The resulting
\href{../results/torsion_phase_transition_checks.json}{\texttt{JSON report}}
is regenerated by running, from the archive's top directory,
\begin{quote}\small\ttfamily
python3 code/verify\_torsion\_phase\_transition.py
\end{quote}
These finite checks supplement the proof; they do not prove its asymptotic assertion.
The argument has also received a separate mathematical reading in this research session,
with the reviewed steps and limitations recorded in
\href{../notes/torsion_review.txt}{\texttt{../notes/torsion\_review.txt}}.
This is not a Lean formalization or external peer review.

To rebuild this document from the archive's top directory, run twice:
\begin{quote}\small\ttfamily\raggedright
pdflatex -interaction=nonstopmode -halt-on-error -output-directory papers
papers/torsion\_phase\_transition.tex
\end{quote}

\begin{thebibliography}{9}\raggedright
\bibitem{proveit-cost}
ProveIt project,
\emph{Torsion specialization: exact costs, obstructions and coefficient fields},
\path{Analysis/GeneralizedRiemannHypothesis/Article/sections/cost.tex},
commit \texttt{b82291c1f75389938c7464ef161907d7d5e75b2a}.
\href{https://github.com/VladimirReshetnikov/ProveIt/blob/b82291c1f75389938c7464ef161907d7d5e75b2a/Analysis/GeneralizedRiemannHypothesis/Article/sections/cost.tex}{Pinned primary source on GitHub}.
Consulted 27 November 2026.

\bibitem{grh-roots}
\texttt{lsernau/grh-lean} contributors,
\emph{RootsOfUnity.lean},
\path{grh/src/New/Congruences/Selmer/RootsOfUnity.lean},
upstream commit \texttt{72fead102f663d5c7a407964a13dc0cf9d977207}.
\href{https://github.com/lsernau/grh-lean/blob/72fead102f663d5c7a407964a13dc0cf9d977207/grh/src/New/Congruences/Selmer/RootsOfUnity.lean}{Pinned upstream primary source on GitHub}.
Vendored in the ProveIt commit above. Consulted 27 November 2026.
\end{thebibliography}
\end{document}
